\chapter{Оршил}

Орчин үед мэдээллийн технологи хурдацтай хөгжиж, нийгмийн бүхий л салбарт гүнзгий нэвтэрч байна. Энэ нөхцөлд тухайн сэдвийг судлах, түүнд суурилсан шийдэл боловсруулах нь онолын хувьд ч, практикийн хувьд ч ихээхэн ач холбогдолтой болоод байна.

Энэхүү дипломын ажлын хүрээнд судалгааны сэдвийн онолын үндсийг авч үзэж, одоо байгаа аргуудыг харьцуулан дүн шинжилгээ хийж, өөрийн санал болгож буй шийдлийг танилцуулна.

\section{Ажлын зорилго}

Дипломын ажлын зорилго нь судалгааны сэдвийн талаарх онолын ойлголтыг системчлэн нэгтгэж, түүнд тулгуурлан практик шийдэл боловсруулж, үр дүнг үнэлэхэд оршино.

Зорилгод хүрэхийн тулд дараах зорилтуудыг дэвшүүлсэн болно:
\begin{itemize}
  \item Сэдэвтэй холбоотой дотоод, гадаадын судалгааны ажлуудыг судлах;
  \item Одоо ашиглагдаж буй аргуудыг харьцуулан дүгнэх;
  \item Шийдлийг боловсруулж, хэрэгжүүлэх;
  \item Үр дүнг туршиж, үнэлэх.
\end{itemize}

\section{Ажлын үндэслэл}

Сэдвийн хамаарал нь одоогийн практикт тулгарч буй асуудал, түүнийг шийдэх хэрэгцээ шаардлагаас үүдэлтэй. Энэ хэсэгт сэдвийг сонгосон үндэслэл, судалгааны шинэлэг тал, практик ач холбогдлыг тодорхойлно.

\section{Ажлын бүтэц}

Дипломын ажил нь оршил, онолын хэсэг, дүгнэлт зэрэг бүлгүүдээс бүрдэнэ. Оршилд ажлын зорилго, үндэслэлийг, онолын хэсэгт судалгааны үндсэн ойлголт, аргуудыг, дүгнэлтэд гарсан үр дүн болон цаашдын чиглэлийг тус тус тусгасан болно.
